\addchap*{Abstract}

Wechselwirkungen zwischen Arzneimitteln sind eine häufige und vermeidbare Ursache unerwünschter Arzneimittelereignisse. Die maßgeblichen Informationen sind zwar öffentlich zugänglich, liegen jedoch überwiegend als fachsprachliche, englische Fachinformationen vor, die für medizinische Laien kaum verständlich sind. Gleichzeitig zeigen aktuelle Studien, dass Large Language Models (LLMs) Wechselwirkungen zwar mit hoher Sensitivität, aber geringer Spezifität erkennen, wenn sie aus ihrem eigenen Trainingswissen antworten.

Diese Arbeit stellt mit \emph{Medcheck} einen Prototyp vor, der beide Beobachtungen verbindet: Das Sprachmodell wird nicht als Wissensquelle, sondern ausschließlich als Erklärwerkzeug über amtliche Daten eingesetzt. Frei eingegebene Medikamentennamen werden über RxNorm normalisiert, der Abschnitt \enquote{Drug Interactions} der Fachinformation wird aus der openFDA Drug Label API bezogen, und ein austauschbares LLM (Anthropic Claude, Google Gemini oder ein lokales Modell über Ollama) überführt diesen Text in eine dreistufige Risikoeinstufung und eine laienverständliche Erklärung in einer von elf Sprachen. Die Implementierung auf Basis von Spring Boot und Angular legt besonderen Wert auf Anbieterunabhängigkeit und auf kontrollierte Degradation bei Ausfall externer Dienste.

In einer Befragung mit elf Teilnehmenden wurde die von Medcheck erzeugte Erklärung zur Kombination von Warfarin und Ibuprofen mit dem entsprechenden Text von Drugs.com verglichen. Die Medcheck-Variante erzielte im Mittel höhere Werte für Verständlichkeit (4{,}30 gegenüber 3{,}64 auf einer fünfstufigen Skala) und wurde von acht der elf Personen als verständlicher bewertet. Beim Vertrauen lagen beide Varianten nahezu gleichauf; mehrere Teilnehmende vermissten in der kurzen Erklärung konkrete Warnzeichen. Die Unterschiede sind bei dieser Stichprobengröße statistisch nicht signifikant. Die Ergebnisse sprechen für eine kombinierte Darstellung aus knapper Zusammenfassung und aufklappbaren Detailinformationen, wie sie der Prototyp bereits vorsieht.
